\documentclass[a4paper]{article}

\usepackage[fontset=fandol]{ctex}
\usepackage{amsmath, amssymb}
\usepackage{graphicx}
\usepackage[margin=2.5cm]{geometry}
\usepackage[most]{tcolorbox}
\usepackage{listings}
\usepackage{hyperref}
\usepackage{booktabs}
\usepackage{float}
\usepackage{tikz}

\setlength{\parindent}{2em}
\setlength{\parskip}{0.35em}
\sloppy

\newtcolorbox{knowledgebox}[1]{
    enhanced, colback=blue!5!white, colframe=blue!75!black, colbacktitle=blue!75!black,
    coltitle=white, fonttitle=\bfseries, title={#1}, attach boxed title to top left={yshift=-2mm, xshift=2mm},
    boxrule=1pt, sharp corners
}

\newtcolorbox{importantbox}[1]{
    enhanced, colback=yellow!10!white, colframe=yellow!80!black, colbacktitle=yellow!80!black,
    coltitle=black, fonttitle=\bfseries, title={#1}, sharp corners
}

\newtcolorbox{warningbox}[1]{
    enhanced, colback=red!5!white, colframe=red!75!black, colbacktitle=red!75!black,
    coltitle=white, fonttitle=\bfseries, title={#1}, sharp corners
}

\lstset{
    basicstyle=\ttfamily\small,
    keywordstyle=\color{blue},
    stringstyle=\color{red!60!black},
    commentstyle=\color{green!60!black},
    breaklines=true,
    frame=single,
    numbers=left,
    numberstyle=\tiny\color{gray},
    captionpos=b,
    extendedchars=false
}

\newcommand{\notetitle}{Transformer（上）：Seq2seq 与 Encoder}
\newcommand{\notesubtitle}{李宏毅《机器学习》2025 第六节：从序列到序列模型到 Transformer 编码器}
\newcommand{\noteauthors}{基于李宏毅老师课程整理}
\newcommand{\notedate}{\today}
\newcommand{\videochannel}{李宏毅机器学习深度学习课程（Bilibili）}
\newcommand{\videoduration}{约 33 分钟}
\newcommand{\videourl}{https://www.bilibili.com/video/BV1TAtwzTE1S?p=15}

\begin{document}
\pagestyle{empty}

\begin{center}
    \vspace*{1.5cm}
    {\Huge\bfseries \notetitle\par}
    \vspace{0.4cm}
    {\LARGE \notesubtitle\par}
    \vspace{0.8cm}
    \includegraphics[width=0.62\textwidth]{video.jpg}\par
    \vspace{0.8cm}
    {\large \noteauthors\par}
    {\large \notedate\par}
    \vspace{0.8cm}
    \begin{tcolorbox}[width=0.82\textwidth,center,sharp corners,
                      colback=black!3!white,colframe=black!50!white]
        \centering
        \textbf{视频信息}\\[3pt]
        频道：\videochannel \\
        时长：\videoduration \\
        链接：\href{\videourl}{\videourl}
    \end{tcolorbox}
\end{center}

\newpage
\tableofcontents
\newpage
\pagestyle{plain}

% =====================================================================
\section{Transformer 从 Seq2seq 开始}
% =====================================================================

这讲课的入口不是一上来写 Transformer 的公式，而是先问一个更基础的问题：什么样的任务需要模型“读入一个序列，再产生另一个序列”？Transformer 本质上就是一个 sequence-to-sequence model，通常简称为 Seq2seq。这里的 sequence 不一定是文字序列，也可以是声音讯号、影像切出来的 patch、标签序列、树结构展开后的符号序列，甚至是模型自己逐步生成的一串动作。

Seq2seq 的关键不是“输入和输出都是序列”这么表面，而是输出长度通常不能由输入长度直接决定。语音辨识是最典型的例子：输入是一段长度为 $T$ 的声音讯号，输出是一串中文字或 token，长度为 $N$。$T$ 和 $N$ 当然有关，但不能只靠 $T$ 推出 $N$；模型必须听懂内容以后，自己决定要输出几个 token。

\begin{figure}[H]
    \centering
    \includegraphics[width=0.88\textwidth]{figures/fig01_seq2seq_speech.jpg}
    \caption{Seq2seq 的基本形式：输入是一段声音序列，输出是一段文字序列，输出长度由模型根据内容决定。\protect\footnotemark}
    \label{fig:seq2seq-speech}
\end{figure}
\footnotetext{视频画面时间区间：约 02:10--02:30；字幕对应“输入的声音讯号它的长度是大 T”“并没有办法一知道输出长度 N 一定是多少”“由机器自己决定”。}

如果把输入序列记作 $X=(x_1,\ldots,x_T)$，输出序列记作 $Y=(y_1,\ldots,y_N)$，Seq2seq 学到的是条件分布
\[
    P(Y\mid X).
\]
在自回归生成的 decoder 中，这个分布通常会拆成
\[
    P(Y\mid X)
    =
    \prod_{t=1}^{N} P(y_t \mid y_{<t}, X).
\]
其中：
\begin{itemize}
    \item $X$ 是输入序列，可以是声音 frame、文字 token、影像 patch 等；
    \item $Y$ 是输出序列，长度 $N$ 往往由模型生成到结束符号时决定；
    \item $y_{<t}$ 表示已经产生的前缀，模型生成第 $t$ 个 token 时会参考前面已经生成的内容。
\end{itemize}

\begin{importantbox}{Seq2seq 的核心能力}
    Seq2seq 不是把每个输入位置逐一对应到一个输出位置，而是把完整输入序列当作条件，再一步一步产生输出序列。它适合处理输入输出长度不固定、对齐关系不明显、输出还需要符合语言或结构约束的任务。
\end{importantbox}

\subsection{和前面 Self-attention 的关系}

前两讲已经介绍过 self-attention：它擅长把一个向量集合或向量序列转换成另一组向量表示。Transformer 把 self-attention 放进 encoder 和 decoder 里面，形成更大的 Seq2seq 系统。也就是说，self-attention 是 Transformer 的局部计算机制；Seq2seq 则是 Transformer 要解决的整体任务形态。

这个区分很重要。只会写 attention 公式，不等于理解 Transformer；理解 Transformer，也不能只停在“输入 token 之间互相看”。真正的模型结构还要回答：如何编码整段输入？如何逐步产生输出？如何把位置信息、残差连接、归一化和全连接网络组合成可训练的深层架构？

\subsection{本章小结}

Seq2seq 把机器学习任务从“固定维度输入到固定类别或数值”推广到“序列输入到序列输出”。Transformer 是 Seq2seq 模型的一种强大实现，而它内部的主要构件就是 self-attention、残差连接、LayerNorm、Feed Forward 和位置编码。

% =====================================================================
\section{为什么许多任务都能被看成 Seq2seq}
% =====================================================================

老师用大量例子说明：Seq2seq 像一把瑞士刀，很多表面上不像“翻译”的任务，只要能把输入和输出都写成序列，就可以硬套成 Seq2seq。这个观点的价值不在于鼓励所有任务都用同一个模型暴力解决，而在于帮助我们看到“输出格式”本身也可以由模型学习。

\subsection{语音翻译：绕开中间文字的瓶颈}

语音辨识通常先把声音转成文字，再让人或另一个模型理解文字。但如果某种语言的文字系统不普及，或者使用者看不懂转写出来的字，单纯 ASR 并不能完成沟通任务。闽南语例子里，机器把台语声音辨识成“母汤”这类文字，听得懂台语的人可能知道意思，但一般使用者未必能理解；如果模型能直接输出普通话“不可”“不行”，任务目标就更接近真实需求。

\begin{figure}[H]
    \centering
    \includegraphics[width=0.86\textwidth]{figures/fig02_hokkien_translation.jpg}
    \caption{同一段闽南语声音可以走“语音辨识”或“语音翻译”两条路径；后者直接输出目标语言，避免把不普及的书写形式当作最终结果。\protect\footnotemark}
    \label{fig:hokkien}
\end{figure}
\footnotetext{视频画面时间区间：约 05:15--05:55；字幕对应“台语的文字并不是一般人能够看得懂”“直接输出的是同样意思的中文的句子”。}

这类任务凸显了 Seq2seq 的一个实用优点：中间表示不一定要是人类可读的文字。模型可以从声音直接到目标语言文字，也可以从文字直接到语音，甚至从一种模态直接到另一种模态。只要训练资料提供输入输出配对，Seq2seq 就有机会学习这个映射。

\begin{warningbox}{端到端不是总是更好}
    直接 speech translation 可以减少错误传播，但也会吃掉中间监督信号。如果 ASR 资料很多、翻译资料很少，分阶段系统可能更容易训练；如果中间文字本身不可靠或不适合用户，端到端模型反而更自然。模型设计要看资料、任务目标和错误容忍度。
\end{warningbox}

\subsection{问答统一视角：把 NLP 任务改写成 QA}

很多自然语言处理任务也可以被包装成问答。给模型一段 context，再给一个 question，模型产生 answer。看似不同的任务，例如翻译、摘要、自然语言推理、情感分析，都可以改写成类似格式：问题描述任务，背景文本提供条件，答案就是任务输出。

\begin{figure}[H]
    \centering
    \includegraphics[width=0.88\textwidth]{figures/fig03_qa_as_seq2seq.jpg}
    \caption{decaNLP 的例子展示了 QA 统一视角：翻译、摘要、推理、情感分析等任务都能写成 question-context-answer 的序列到序列格式。\protect\footnotemark}
    \label{fig:qa}
\end{figure}
\footnotetext{视频画面时间区间：约 12:05--12:35；字幕对应“很多 natural language processing 的任务都可以想成是 question answering 的任务”。}

这个统一视角对现代大模型尤其重要。今天的 instruction following、prompt-based learning、text-to-text framework，背后都有类似思想：把任务说明和输入资料放进一个序列，再让模型生成目标序列。任务的边界从模型结构转移到输入输出格式上，模型本身可以保持相对统一。

\begin{knowledgebox}{统一格式的好处}
    当多个任务都被写成“读入一串 token，输出一串 token”，模型可以共享同一套参数和训练流程。不同任务之间的知识可能互相迁移，例如阅读理解训练帮助摘要，翻译训练帮助语义对齐，分类任务则可以被看成只输出很短答案的生成任务。
\end{knowledgebox}

\subsection{树结构也能线性化}

句法分析表面上输出的是一棵树，不是序列。老师的例子是句子 “deep learning is very powerful”：模型需要判断 “deep learning” 是名词片语， “very powerful” 是形容词片语，再和 “is” 组成动词片语，最后得到完整句子结构。

树不是序列，但可以写成括号表达式，例如
\[
    (S\ (NP\ deep\ learning)\ (VP\ is\ (ADJV\ very\ powerful))).
\]
一旦树被线性化，它就变成模型可以逐 token 产生的输出序列。模型不必直接输出一个复杂数据结构，而是输出一种能被解析回树的符号串。

\begin{figure}[H]
    \centering
    \includegraphics[width=0.86\textwidth]{figures/fig04_parsing_as_sequence.jpg}
    \caption{句法分析可以把树状输出线性化成括号序列；Seq2seq 因而能处理看起来不是序列输出的结构预测任务。\protect\footnotemark}
    \label{fig:parsing}
\end{figure}
\footnotetext{视频画面时间区间：约 19:35--19:55；画面展示括号化的 parsing sequence 与相关论文，字幕讨论 “sequence to sequence model 居然可以硬做到”。}

这种技巧的意义不止于 parsing。很多结构化预测都可以被“序列化”：多标签分类可以输出标签序列，目标检测可以输出类别与 bounding box 序列，程序生成可以输出代码 token，数学推理可以输出证明步骤。Seq2seq 的强大之处正在于它把“输出空间的形状”交给 token 序列承载。

\subsection{Seq2seq 不是万能最优解}

老师反复强调，Seq2seq 可以硬做很多事情，但不代表每件事情都应该这样做。任务专用模型常常更高效、更可解释、更容易满足约束。例如目标检测若直接生成 box 序列，需要额外处理重复框、顺序、坐标合法性等问题；专门的 detection architecture 会把这些先验写进模型结构中。

\begin{importantbox}{瑞士刀比喻的正确读法}
    Seq2seq 是一个通用建模语言：它让我们知道“很多任务原则上可以统一”。但工程上仍要比较资料量、速度、准确率、约束满足、错误分析和部署成本。统一模型提供可能性，专门模型提供归纳偏置。
\end{importantbox}

\subsection{本章小结}

语音翻译、问答、摘要、情感分析、句法分析、多标签分类和目标检测都可以用 Seq2seq 描述。这个统一视角让我们从“为每个任务设计一个输出头”转向“把任务输出写成可生成的序列”。但通用性不等于最优性，实际建模仍需要考虑任务结构和资料条件。

% =====================================================================
\section{从传统 Seq2seq 到 Transformer}
% =====================================================================

Seq2seq 模型的历史早于 Transformer。2014 年左右，基于 RNN/LSTM 的 encoder-decoder 已经被用于机器翻译：encoder 读完整个输入句子，decoder 再逐步产生输出句子。那时的模型相对朴素，长距离依赖和并行训练都比较困难。

Transformer 的突破在于，它用 attention 取代 RNN 的主要序列建模工作。输入不再必须按时间步一个接一个递推；每个位置可以通过 self-attention 直接参考其他位置。对长序列、并行训练和大规模资料而言，这个结构更适合现代硬件，也更容易堆深。

\begin{figure}[H]
    \centering
    \includegraphics[width=0.88\textwidth]{figures/fig05_seq2seq_to_transformer.jpg}
    \caption{传统 Seq2seq 和 Transformer 都是 encoder-decoder 架构；Transformer 论文中的彩色模块把 attention、Add \& Norm 和 Feed Forward 组织成可堆叠结构。\protect\footnotemark}
    \label{fig:seq2seq-transformer}
\end{figure}
\footnotetext{视频画面时间区间：约 23:35--24:00；字幕对应“14 年就有 sequence to sequence model 用在翻译”“今天主角也就是 Transformer”“有 encoder 架构，有 decoder 架构”。}

从系统角度看，Seq2seq 可以抽象为
\[
    \text{Encoder}: X \mapsto H,\qquad
    \text{Decoder}: (H, y_{<t}) \mapsto y_t.
\]
其中：
\begin{itemize}
    \item $X=(x_1,\ldots,x_T)$ 是输入序列；
    \item $H=(h_1,\ldots,h_T)$ 是 encoder 输出的上下文化表示；
    \item decoder 在生成每个输出 token 时，会参考 $H$ 和已经生成的输出前缀。
\end{itemize}

本讲后半主要展开 encoder。decoder 的细节会留到下一讲，所以这里先把 encoder 看清楚：它接收一排向量，输出另一排同长度向量；每个输出向量都应该包含对整个输入序列的理解。

\subsection{Encoder 和 BERT 的关系}

老师在开头用 BERT 暗示 Transformer 与后续模型的关系。严格说，BERT 可以理解为主要使用 Transformer encoder 的模型：它不做自回归输出序列生成，而是把输入序列编码成上下文化表示，再用这些表示完成遮蔽词预测、句子关系判断或下游任务微调。

因此，学 Transformer encoder 不只是为了完整 Seq2seq，也是在为理解 BERT、现代 embedding model、文本分类模型和许多视觉 Transformer 打基础。encoder 的问题是“如何得到好的表示”；decoder 的问题是“如何根据表示生成东西”。

\subsection{本章小结}

Transformer 继承了 Seq2seq 的 encoder-decoder 框架，但内部计算从 RNN 递推转向 self-attention 和可并行堆叠的 block。本讲先聚焦 encoder：输入一排向量，输出一排带上下文信息的向量。

% =====================================================================
\section{Transformer Encoder：一排向量到另一排向量}
% =====================================================================

在前面 self-attention 课程里，已经讨论过“输入是一组向量，输出也是一组向量”的问题。Transformer encoder 正是这种问题的标准答案之一。给定输入 token 的 embedding
\[
    x^1,x^2,\ldots,x^T,
\]
encoder 输出
\[
    h^1,h^2,\ldots,h^T.
\]
输出长度和输入长度相同，但每个 $h^i$ 已经不只是 $x^i$ 的局部变换，而是通过 self-attention 吸收了其他位置的信息。

\begin{figure}[H]
    \centering
    \includegraphics[width=0.82\textwidth]{figures/fig06_encoder_stack.jpg}
    \caption{Transformer encoder 由多个 block 堆叠而成；每个 block 都把一排向量变成另一排向量，最后得到上下文化表示 $h^1,\ldots,h^T$。\protect\footnotemark}
    \label{fig:encoder-stack}
\end{figure}
\footnotetext{视频画面时间区间：约 24:45--25:10；字幕对应“encoder 里面会分成很多很多的 block”“每一个 block 都是输入一排向量，输出一排向量”。}

\subsection{为什么不能只写一个 Self-attention}

一个 self-attention 层能让不同位置交换信息，但 Transformer encoder 不是只把 self-attention 单独摆在那里。原因有三点：
\begin{itemize}
    \item self-attention 主要做跨位置的信息混合，但每个位置内部还需要非线性变换；
    \item 深层网络训练需要残差连接来保留梯度通路；
    \item 中间表示的尺度需要归一化，否则堆叠很多层后训练容易不稳定。
\end{itemize}

因此，一个 encoder block 通常由两类子层组成：multi-head self-attention 和 position-wise feed-forward network。每个子层外面再接 residual connection 和 layer normalization。原始 Transformer 论文图里把这些合称为 Add \& Norm。

\begin{knowledgebox}{Position-wise Feed Forward 是什么}
    Feed Forward Network 不是在不同 token 之间交换信息，而是对每个位置的向量各自做同一组全连接变换。跨位置的信息混合主要由 self-attention 完成；每个位置内部的特征重组、升维、非线性和降维主要由 FFN 完成。
\end{knowledgebox}

\subsection{位置编码不可省}

Self-attention 本身对输入顺序不敏感。如果交换输入向量的顺序，attention 机制只看到一组向量，无法天然知道谁在前谁在后。Transformer 因此要把 positional encoding 加到 input embedding 上，让每个位置带有位置信息。

位置编码可以是固定的正弦余弦函数，也可以是学习出来的向量。无论形式如何，它解决的是同一个问题：让模型不只知道 token 内容，也知道 token 在序列中的位置。

\subsection{本章小结}

Transformer encoder 是一个可堆叠的向量序列变换器。它保留输入输出长度一致的形式，但通过多层 self-attention、FFN、残差和归一化，让每个输出向量包含全局上下文。位置编码负责把顺序信息注入原本 permutation-invariant 的 attention 计算。

% =====================================================================
\section{一个 Encoder Block 的内部结构}
% =====================================================================

老师用逐步展开的图解释一个 encoder block。第一步是 self-attention：输入一排向量，输出一排经过上下文混合的向量。接着不是直接把它传给下一层，而是先做 residual connection，也就是把 self-attention 的输入 $a$ 和输出 $b$ 相加：
\[
    u = a + b.
\]
这里 $a$ 是子层输入，$b=\mathrm{SelfAttention}(a)$ 是子层输出，$u$ 是残差相加后的结果。

残差连接的直觉是保留一条“原始信息和梯度都比较容易通过”的路径。深层网络如果每一层都必须完整重写表示，训练会困难；残差让每层更像是在学习一个修正量，而不是从零开始生成新表示。

\subsection{Layer Normalization}

Residual 之后接 LayerNorm。对一个向量
\[
    x=(x_1,x_2,\ldots,x_K),
\]
LayerNorm 会在同一个样本、同一个 token 的不同维度上计算均值和标准差：
\[
    m=\frac{1}{K}\sum_{i=1}^{K}x_i,
    \qquad
    \sigma=\sqrt{\frac{1}{K}\sum_{i=1}^{K}(x_i-m)^2+\epsilon},
\]
\[
    x_i'=\frac{x_i-m}{\sigma}.
\]
其中：
\begin{itemize}
    \item $K$ 是 hidden dimension；
    \item $m$ 和 $\sigma$ 都来自同一个向量内部的各个维度；
    \item $\epsilon$ 是避免除以零的小常数；
    \item 实际实现通常还会有可学习的缩放和平移参数 $\gamma,\beta$，即 $\gamma x_i'+\beta$。
\end{itemize}

\begin{figure}[H]
    \centering
    \includegraphics[width=0.88\textwidth]{figures/fig07_encoder_block_residual_layernorm_ffn.jpg}
    \caption{一个 encoder block 中，self-attention 后接 residual 与 LayerNorm，之后进入 FC/Feed Forward，并再次使用 residual 结构。\protect\footnotemark}
    \label{fig:block}
\end{figure}
\footnotetext{视频画面时间区间：约 28:25--29:05；字幕对应“每一个 dimension 减掉 mean 再除以 standard deviation”“fully connected network 这边也有 residual 的架构”。}

\begin{warningbox}{LayerNorm 和 BatchNorm 的统计方向不同}
    BatchNorm 通常对同一个 feature dimension、不同 example 计算统计量；LayerNorm 则对同一个 example 内部的不同 hidden dimensions 计算统计量。Transformer 常用 LayerNorm，是因为序列长度、batch 组成和自回归生成场景会让 batch-level statistics 更不稳定。
\end{warningbox}

\subsection{Feed Forward Network 的角色}

经过 self-attention 后，每个 token 已经获得其他 token 的信息，但这个混合后的向量还需要进一步非线性加工。FFN 通常写成
\[
    \mathrm{FFN}(x)=W_2\,\phi(W_1x+b_1)+b_2,
\]
其中 $\phi$ 可以是 ReLU、GELU 等非线性函数。这个 FFN 对每个位置独立套用同一组参数，所以它不会改变序列长度。

一个常见 encoder block 可以抽象写成：
\[
    z = \mathrm{LayerNorm}(x + \mathrm{MHA}(x)),
\]
\[
    h = \mathrm{LayerNorm}(z + \mathrm{FFN}(z)).
\]
这里 $\mathrm{MHA}$ 是 multi-head attention，$z$ 是 attention 子层后的表示，$h$ 是整个 block 的输出。

需要注意：上面写法对应原始 Transformer 常见的 post-norm 形式。后来很多模型会改成 pre-norm，也就是先 LayerNorm 再进入子层。两者训练稳定性不同，后面课程和论文会进一步讨论。

\subsection{本章小结}

Encoder block 的内部并不只是 self-attention。一个完整 block 至少包含 self-attention、residual connection、LayerNorm、FFN，以及 FFN 外面的 residual/normalization。Self-attention 负责跨位置沟通，FFN 负责位置内非线性变换，residual 和 LayerNorm 负责让深层堆叠可训练。

% =====================================================================
\section{对照原始 Transformer Encoder 图}
% =====================================================================

课程最后把刚才拆开的结构对回原始 Transformer 论文的 encoder 图。原图右侧可以看到：输入先经过 input embedding，再加 positional encoding；之后进入 $N$ 次重复的 encoder block。每个 block 内部有两个大子层：Multi-Head Attention 和 Feed Forward。每个子层旁边都有 Add \& Norm。

\begin{figure}[H]
    \centering
    \includegraphics[width=0.88\textwidth]{figures/fig08_transformer_encoder_original.jpg}
    \caption{原始 Transformer encoder 图中，Add \& Norm 表示 residual connection 加 LayerNorm；Multi-Head Attention 对应 self-attention 子层，Feed Forward 对每个位置独立变换。\protect\footnotemark}
    \label{fig:original-encoder}
\end{figure}
\footnotetext{视频画面时间区间：约 29:40--30:25；字幕对应“positional information”“multi-head self attention”“Add and Norm 就是 residual 加 layer normalization”。}

这张图里有几个容易漏掉的细节：
\begin{itemize}
    \item 输入 embedding 和 positional encoding 是相加，不是串接；
    \item encoder block 重复 $N$ 次，所以深度来自同一结构的堆叠；
    \item Multi-Head Attention 是 self-attention 的多头版本，每个 head 可以关注不同关系；
    \item Add \& Norm 出现两次，分别包住 attention 子层和 FFN 子层；
    \item encoder 输出是一排向量，不是一个单一向量。
\end{itemize}

\subsection{Multi-head 的意义}

单头 attention 会把每个 token 对其他 token 的关系压进一组 attention weights。Multi-head attention 则把表示分成多个子空间，让不同 head 可以学不同关系：有的 head 可能关注局部搭配，有的关注长距离依赖，有的关注语法角色，有的关注实体对应。虽然不应把每个 head 都机械解释成某种固定语言学功能，但多头结构确实增加了模型同时表达多种关系的能力。

\subsection{为什么这里讲 LayerNorm 而不是 BatchNorm}

前一讲刚讲 Batch Normalization，所以这讲特别对比 LayerNorm。BatchNorm 的统计量依赖 batch 内其他样本；LayerNorm 的统计量来自单个样本的 hidden dimensions。对 Transformer 来说，LayerNorm 有几个优势：
\begin{itemize}
    \item 推论时不需要维护 batch-level running mean/variance；
    \item 适合 batch size 变化较大的训练与推论场景；
    \item 自回归生成时一次可能只处理少量 token，BatchNorm 的 batch 统计不自然；
    \item 序列任务中每个位置的表示更适合在 hidden dimension 内稳定尺度。
\end{itemize}

这不是说 BatchNorm 不能用于 Transformer，而是原始架构和许多后续模型更偏向 LayerNorm。后来也有研究重新审视 normalization 的位置和形式，说明 Transformer 的细节设计并非不可改变。

\subsection{本章小结}

原始 Transformer encoder 图可以拆成清楚的四件事：embedding 加位置、multi-head self-attention、position-wise FFN、两处 Add \& Norm。看懂这张图，就能把前面讲过的 self-attention 和后续会讲的 BERT、decoder、Seq2seq 训练连接起来。

% =====================================================================
\section{Normalization 位置不是唯一答案}
% =====================================================================

课程结尾提到：原始论文的结构不代表永远最优。比如 LayerNorm 应该放在 residual 之后，还是放在每个子层之前？这就是 post-norm 和 pre-norm 的差异。原始 Transformer 图更接近 post-norm；许多后来大模型为了训练更深网络，会采用 pre-norm 或其他变体。

\begin{figure}[H]
    \centering
    \includegraphics[width=0.86\textwidth]{figures/fig09_layernorm_design_variants.jpg}
    \caption{课程给出的拓展阅读强调：Transformer 中 LayerNorm 的位置和 normalization 形式都可以研究，原始设计不是唯一可行结构。\protect\footnotemark}
    \label{fig:ln-variants}
\end{figure}
\footnotetext{视频画面时间区间：约 31:15--32:20；字幕对应“为什么 layer normalization 是放在那个地方呢”“原始论文的设计并不代表它是最好的”。}

\subsection{Pre-norm 与 Post-norm 的直觉差异}

用简化形式表示，post-norm 可以写成
\[
    x_{l+1}=\mathrm{LN}(x_l + F(x_l)).
\]
pre-norm 则可以写成
\[
    x_{l+1}=x_l + F(\mathrm{LN}(x_l)).
\]
两者都包含 residual 和 LayerNorm，但梯度路径不同。pre-norm 中，残差主路径更直接，深层模型训练通常更稳定；post-norm 中，归一化直接作用在相加后的输出上，原始 Transformer 使用这种形式，但堆很深时可能更难训练。

\begin{knowledgebox}{结构细节会影响可扩展性}
    Transformer 的大方向是 attention + FFN + residual + normalization，但真正训练大模型时，LayerNorm 位置、激活函数、初始化、残差缩放、学习率计划都会影响稳定性。架构图看起来只是几个方块，工程上每个方块的放置都可能决定模型能否顺利扩大。
\end{knowledgebox}

\subsection{本章小结}

原始 Transformer encoder 是理解后续模型的基准，但不是唯一答案。Normalization 的位置和形式仍然是可研究、可替换的设计空间。学架构时应该先掌握标准版，再知道哪些地方是历史选择、工程折衷或可调超参数。

% =====================================================================
\section{总结与延伸}
% =====================================================================

本讲完成了两个层次的铺垫。第一层是任务层：Transformer 是 Seq2seq model，而 Seq2seq 可以把语音辨识、语音翻译、问答、摘要、句法分析、多标签分类和目标检测等任务放到同一个“输入序列到输出序列”的框架下。第二层是结构层：Transformer encoder 把一排输入向量变成一排上下文化向量，内部由 multi-head self-attention、FFN、residual connection、LayerNorm 和 positional encoding 组成。

可以把本讲压缩成下面这条链：
\[
\begin{aligned}
    \text{Seq2seq task}
    &\rightarrow
    \text{Encoder-Decoder framework}
    \rightarrow
    \text{Transformer}\\
    &\rightarrow
    \text{Self-attention blocks}
    \rightarrow
    \text{Residual + LayerNorm + FFN}.
\end{aligned}
\]
其中每一步都解决一个具体问题：
\begin{itemize}
    \item Seq2seq 解决输入输出长度不固定的问题；
    \item encoder-decoder 分工解决“先理解输入，再产生输出”的问题；
    \item self-attention 解决序列中远距离位置互相参考的问题；
    \item multi-head 解决关系类型多样的问题；
    \item positional encoding 解决顺序信息缺失的问题；
    \item residual 和 LayerNorm 解决深层网络训练稳定性的问题；
    \item FFN 解决每个位置内部非线性加工的问题。
\end{itemize}

\begin{importantbox}{本讲最重要的理解}
    Transformer encoder 不是一个神秘黑盒，而是“向量序列变换”的可堆叠模块。它每层先让 token 之间互相交换信息，再对每个 token 的表示做非线性加工，并用残差与归一化维持训练稳定。这个结构既能作为完整 Seq2seq 的 encoder，也能单独成为 BERT 这类表示学习模型的核心。
\end{importantbox}

下一讲自然会进入 decoder。Encoder 负责把输入序列变成表示；decoder 则要回答另一个难题：输出序列是一步一步生成的，模型在产生第 $t$ 个 token 时不能偷看未来 token，同时还要能参考 encoder 的输出。于是 masked self-attention、cross-attention、teacher forcing、结束符号等机制都会登场。

\subsection{拓展阅读}

本讲画面中提到的几篇资料可以作为延伸：
\begin{itemize}
    \item \href{https://arxiv.org/abs/1706.03762}{Attention Is All You Need}：原始 Transformer 架构；
    \item \href{https://arxiv.org/abs/1412.7449}{Grammar as a Foreign Language}：把 parsing 视为序列生成；
    \item \href{https://arxiv.org/abs/2002.04745}{On Layer Normalization in the Transformer Architecture}：讨论 LayerNorm 位置；
    \item \href{https://arxiv.org/abs/2003.07845}{PowerNorm}：重新思考 Transformer 中的 normalization。
\end{itemize}

\end{document}
